%% preamble for the handout

\documentclass[titlepage,twoside,parskip=half,BCOR14mm,DIV13]{scrartcl}
\usepackage[T1]{fontenc}
\usepackage[utf8]{inputenc}
\usepackage[absolute]{textpos}
\usepackage{beamerarticle}
\usepackage{beamerthemeshadow}
\usepackage{hyperref}
\usepackage{longtable,makecell}
\usepackage{multicol}
\usepackage{tabularx}
\usepackage{multirow}
\usepackage{pgf}
\usepackage{float}
\usepackage{xspace}
\usepackage{rotating}
\usepackage{ulem}
\usepackage{courier}
\usepackage{mathpazo}
\usepackage[scaled=.95]{helvet}


\usepackage{scrlayer-scrpage}
\usepackage{textcomp}
%\usepackage{titleref} %% Use this if you need to reference to a section.
\usepackage{typearea}
\usepackage[tikz]{bclogo}

\pagestyle{scrheadings}
\automark[section]{section}
\renewcommand{\headfont}{\upshape}
\ofoot{\thepage}
\cfoot{}

\beamersetuncovermixins{\opaqueness<1>{25}}{\opaqueness<2->{15}}

\usepackage{listings}
\lstset{
  breaklines=true,
  breakatwhitespace=true,
  postbreak=\raisebox{0ex}[0ex][0ex]{\ensuremath{\hookrightarrow\space}},
  basicstyle=\small\ttfamily,
  extendedchars=true,
  inputencoding=utf8,
  frame=none,
  literate={Ö}{{\"O}}1 {Ä}{{\"A}}1 {Ü}{{\"U}}1 {ß}{{\ss}}2 {ü}{{\"u}}1 {ä}{{\"a}}1 {ö}{{\"o}}1 {µ}{\textmu}1
}
\usepackage[\configLang]{babel}
\usepackage{graphicx}
\usepackage{color}
\usepackage[gen]{eurosym}

%\usepackage{pgf}

% Wasserzeichen im PDF
%\usepackage{draftwatermark}
%\SetWatermarkText{Internal use only}
%\SetWatermarkFontSize{3cm}

% Paket für Profiling
\RequirePackage{versions}



%\setjobnamebeamerversion{lvm}
\usetheme{Luebeck}
\useoutertheme{infolines}

\definecolor{b1blau}{RGB}{34,71,121}
\definecolor{b1orange}{RGB}{252,189,43}
\definecolor{grau}{RGB}{40,40,40}
\definecolor{greylist}{RGB}{224,224,224}
\definecolor{tmagenta}{RGB}{226,0,116}
\setbeamercolor{structure}{fg=b1blau}
\setbeamercolor*{palette primary}{bg=b1blau}
\setbeamercolor*{palette secondary}{bg=b1blau}
\setbeamercolor*{palette tertiary}{bg=b1blau}

\setbeamercolor{palette primary}{use=structure,fg=white,bg=b1blau!100!black}
\setbeamercolor{palette secondary}{use=structure,fg=white,bg=b1blau!100!white}
\setbeamercolor{palette tertiary}{use=structure,fg=white,bg=b1blau!100!white}
\setbeamercolor{palette quaternary}{use=structure,fg=white,bg=b1blau!100!white}

% Definitionen für Hinweisgrafiken
\renewcommand{\bcdanger}{\includegraphics[width=0.85\logowidth]{common/artwork/bc_danger_icon}}
\renewcommand{\bcattention}{\includegraphics[width=0.85\logowidth]{common/artwork/bc_attention_icon}}
\renewcommand{\bclampe}{\includegraphics[width=0.85\logowidth]{common/artwork/bc_lampe_icon}}
\renewcommand{\bcinfo}{\includegraphics[width=0.85\logowidth]{common/artwork/bc_info_icon}}

% Definitionen für Hinweisboxen
%\newenvironment{name}[num]{before}{after}
\def\boxesgerman{%
 \def\warning{Warnung}%
 \def\important{Wichtig}%
 \def\tip{Tipp}%
 \def\note{Hinweis}%
}
\def\boxesenglish{%
 \def\warning{Warning}%
 \def\important{Important}%
 \def\tip{Tip}%
 \def\note{Note}%
}
\ifthenelse{\equal{\configLang}{english}}{\boxesenglish}{\boxesgerman}

\newenvironment{b1warning}[1]{%
\begin{bclogo}[logo=\bcdanger,couleur=blue!10,barre=none,epBord=0.4,ombre=true,couleurOmbre=black!60,blur,arrondi=0.1,sousTitre=#1]{\warning}}{\end{bclogo}}

\newenvironment{b1important}[1]{%
\begin{bclogo}[logo=\bcattention,couleur=blue!10,barre=none,epBord=0.4,ombre=true,couleurOmbre=black!60,blur,arrondi=0.1,sousTitre=#1]{\important}}{\end{bclogo}}

\newenvironment{b1tip}[1]{%
\begin{bclogo}[logo=\bclampe,couleur=blue!10,barre=none,epBord=0.4,ombre=true,couleurOmbre=black!60,blur,arrondi=0.1,sousTitre=#1]{\tip}}{\end{bclogo}}

\newenvironment{b1note}[1]{%
\begin{bclogo}[logo=\bcinfo,couleur=blue!10,barre=none,epBord=0.4,ombre=true,couleurOmbre=black!60,blur,arrondi=0.1,sousTitre=#1]{\note}}{\end{bclogo}}

\setcounter{section}{-1}
